\begin{enumerate}
    \item Из основания $D$ высоты $AD$ остроугольного треугольника $ABC$ опущены перпендикуляры $DK$ и $DL$ на стороны $AB$ и $AC$ соответственно. Известно, что $\angle BAC = 72^\circ, \angle ABL = 30^\circ$. Чему равен угол $\angle HKC$?

    \item Касательные к описанной окружности остроугольного треугольника $ABC$, проведённые в точках $B$ и $C$, пересекаются в точке $D$. Серединный перпендикуляр к стороне $AC$ пересекает отрезок $AB$ в точке $E$. Найдите угол $BDE$, если $\angle A = 34^\circ, \angle B = 70^\circ$.
    
    \item В треугольнике $ABC$ точки $B_1, C_1$ -- основания высот из вершин $B, C$ соответственно. Точка $D$ -- основания перпендикуляра из $B_1$ на $AB$, $E$ -- точка пересечния перпендикуляра, опущенного из точки $D$ на сторон $BC$, с отрезком $BB_1$. Докажите, что $EC_1 \perp BB_1$.

    \item На гипотенузе $AC$ прямоугольного треугольника $ABC$ во внешнюю сторону треугольника построен квадрат с центром в точке $O$. Докажите, что $BO$ — биссектриса угла $ABC$.

    \item В остроугольном треугольнике $ABC$ на высоте, проведённой из вершины $C$, выбрана точка $P$. Пусть $A_1$ и $B_1$ — основания перпендикуляров из точки $X$ на стороны $AC$ и $BC$ соответственно. Докажите, что точки $A$, $B$, $B_1$, $A_1$ концикличны.

    \item Пусть $AA_1$, $BB_1$, $CC_1$ — высоты остроугольного треугольника $ABC$. Докажите, что основания перпендикуляров из точки $A_1$ на прямые $AB$, $AC$, $BB1$, $CC1$ коллинеарны.

    \item Точки $A$, $B$, $C$ лежат на окружности, а прямая $\ell$ касается этой окружности в точке $B$. Из точки $P$, лежащей на прямой $\ell$, опущены перпендикуляры $PA_1$ и $PC_1$ на прямые $AB, BC$ соответственно. Докажите, что $A_1C_1 \perp AC$.
\end{enumerate}